% ======================================================================
%  UPPGIFT 2: FÖRÄNDRAT FLÖDE
% ======================================================================
\section{2 Förändrat flöde}

\uppgiftsbild{2}{Förändrat flöde}
{uppdatera ett maximalt flöde när en kapacitet ändras med 1 — i linjär tid, utan att börja om.}
{\begin{itemize}\setlength{\itemsep}{1pt}
  \item Vad lösningen $\Phi$ \emph{är}: flödet på varje kant, inte bara värdet.
  \item Riktningen: ökning (högst $+1$) mot minskning (högst $-1$).
  \item Kanten: mättad eller inte — och ligger den i \emph{alla} eller i \emph{något} minsnitt?
  \item ”Linjär tid”: linjär i \emph{vad}?
\end{itemize}}

\begin{frame}{Repetition: Ford--Fulkerson på exempelnätet}
\begin{columns}[T]
\begin{column}{0.49\textwidth}
\centering
{\small\bfseries\color{kthnavy}Flöde $f$}\ {\scriptsize(flöde/kapacitet, \textcolor{brick}{röd} = mättad)}\par\smallskip
\only<1>{%
\begin{tikzpicture}[scale=0.95]
\useasboundingbox (-0.45,-0.5) rectangle (5.85,3.3);
  \node[fn] (s) at (0,1.4) {$s$};
  \node[fn] (a) at (2.6,2.8) {$a$};
  \node[fn] (b) at (1.6,0) {$b$};
  \node[fn] (c) at (3.9,0) {$c$};
  \node[fn] (t) at (5.4,1.4) {$t$};
  \draw[fe] (s) -- node[fl] {0/3} (a);
  \draw[fe] (s) -- node[fl] {0/6} (b);
  \draw[fe] (b) -- node[fl] {0/1} (a);
  \draw[fe] (a) -- node[fl] {0/1} (t);
  \draw[fe] (a) -- node[fl] {0/2} (c);
  \draw[fe] (b) -- node[fl] {0/5} (c);
  \draw[fe] (c) -- node[fl] {0/7} (t);
\end{tikzpicture}}
\only<2>{%
\begin{tikzpicture}[scale=0.95]
\useasboundingbox (-0.45,-0.5) rectangle (5.85,3.3);
  \node[fn] (s) at (0,1.4) {$s$};
  \node[fn] (a) at (2.6,2.8) {$a$};
  \node[fn] (b) at (1.6,0) {$b$};
  \node[fn] (c) at (3.9,0) {$c$};
  \node[fn] (t) at (5.4,1.4) {$t$};
  \draw[fe] (s) -- node[fl] {2/3} (a);
  \draw[fe] (s) -- node[fl] {0/6} (b);
  \draw[fe] (b) -- node[fl] {0/1} (a);
  \draw[fe] (a) -- node[fl] {0/1} (t);
  \draw[sat] (a) -- node[fl] {2/2} (c);
  \draw[fe] (b) -- node[fl] {0/5} (c);
  \draw[fe] (c) -- node[fl] {2/7} (t);
\end{tikzpicture}}
\only<3>{%
\begin{tikzpicture}[scale=0.95]
\useasboundingbox (-0.45,-0.5) rectangle (5.85,3.3);
  \node[fn] (s) at (0,1.4) {$s$};
  \node[fn] (a) at (2.6,2.8) {$a$};
  \node[fn] (b) at (1.6,0) {$b$};
  \node[fn] (c) at (3.9,0) {$c$};
  \node[fn] (t) at (5.4,1.4) {$t$};
  \draw[fe] (s) -- node[fl] {2/3} (a);
  \draw[fe] (s) -- node[fl] {1/6} (b);
  \draw[fe] (b) -- node[fl] {0/1} (a);
  \draw[sat] (a) -- node[fl] {1/1} (t);
  \draw[fe] (a) -- node[fl] {1/2} (c);
  \draw[fe] (b) -- node[fl] {1/5} (c);
  \draw[fe] (c) -- node[fl] {2/7} (t);
\end{tikzpicture}}
\only<4>{%
\begin{tikzpicture}[scale=0.95]
\useasboundingbox (-0.45,-0.5) rectangle (5.85,3.3);
  \node[fn] (s) at (0,1.4) {$s$};
  \node[fn] (a) at (2.6,2.8) {$a$};
  \node[fn] (b) at (1.6,0) {$b$};
  \node[fn] (c) at (3.9,0) {$c$};
  \node[fn] (t) at (5.4,1.4) {$t$};
  \draw[fe] (s) -- node[fl] {2/3} (a);
  \draw[fe] (s) -- node[fl] {5/6} (b);
  \draw[fe] (b) -- node[fl] {0/1} (a);
  \draw[sat] (a) -- node[fl] {1/1} (t);
  \draw[fe] (a) -- node[fl] {1/2} (c);
  \draw[sat] (b) -- node[fl] {5/5} (c);
  \draw[fe] (c) -- node[fl] {6/7} (t);
\end{tikzpicture}}
\only<5>{%
\begin{tikzpicture}[scale=0.95]
\useasboundingbox (-0.45,-0.5) rectangle (5.85,3.3);
  \node[fn] (s) at (0,1.4) {$s$};
  \node[fn] (a) at (2.6,2.8) {$a$};
  \node[fn] (b) at (1.6,0) {$b$};
  \node[fn] (c) at (3.9,0) {$c$};
  \node[fn] (t) at (5.4,1.4) {$t$};
  \draw[fe] (s) -- node[fl] {2/3} (a);
  \draw[sat] (s) -- node[fl] {6/6} (b);
  \draw[sat] (b) -- node[fl] {1/1} (a);
  \draw[sat] (a) -- node[fl] {1/1} (t);
  \draw[sat] (a) -- node[fl] {2/2} (c);
  \draw[sat] (b) -- node[fl] {5/5} (c);
  \draw[sat] (c) -- node[fl] {7/7} (t);
\end{tikzpicture}}
\end{column}
\begin{column}{0.49\textwidth}
\centering
{\small\bfseries\color{kthnavy}Restflödesgraf $G_f$}\ {\scriptsize(\textcolor{kthblue}{kvar}, \textcolor{kthorange}{ångra})}\par\smallskip
\only<1>{%
\begin{tikzpicture}[scale=0.95]
\useasboundingbox (-0.45,-0.5) rectangle (5.85,3.3);
  \node[fn] (s) at (0,1.4) {$s$};
  \node[fn] (a) at (2.6,2.8) {$a$};
  \node[fn] (b) at (1.6,0) {$b$};
  \node[fn] (c) at (3.9,0) {$c$};
  \node[fn] (t) at (5.4,1.4) {$t$};
  \draw[aug] (s) -- node[fl,text=kthgrass,font=\scriptsize\bfseries] {3} (a);
  \draw[res] (s) -- node[fl] {6} (b);
  \draw[res] (b) -- node[fl] {1} (a);
  \draw[res] (a) -- node[fl] {1} (t);
  \draw[aug] (a) -- node[fl,text=kthgrass,font=\scriptsize\bfseries] {2} (c);
  \draw[res] (b) -- node[fl] {5} (c);
  \draw[aug] (c) -- node[fl,text=kthgrass,font=\scriptsize\bfseries] {7} (t);
\end{tikzpicture}}
\only<2>{%
\begin{tikzpicture}[scale=0.95]
\useasboundingbox (-0.45,-0.5) rectangle (5.85,3.3);
  \node[fn] (s) at (0,1.4) {$s$};
  \node[fn] (a) at (2.6,2.8) {$a$};
  \node[fn] (b) at (1.6,0) {$b$};
  \node[fn] (c) at (3.9,0) {$c$};
  \node[fn] (t) at (5.4,1.4) {$t$};
  \draw[res] (s) to[bend left=20] node[fl] {1} (a);
  \draw[resb] (a) to[bend left=20] node[fl] {2} (s);
  \draw[aug] (s) -- node[fl,text=kthgrass,font=\scriptsize\bfseries] {6} (b);
  \draw[res] (b) -- node[fl] {1} (a);
  \draw[aug] (a) -- node[fl,text=kthgrass,font=\scriptsize\bfseries] {1} (t);
  \draw[aug] (c) -- node[fl,text=kthgrass,font=\scriptsize\bfseries] {2} (a);
  \draw[aug] (b) -- node[fl,text=kthgrass,font=\scriptsize\bfseries] {5} (c);
  \draw[res] (c) to[bend left=20] node[fl] {5} (t);
  \draw[resb] (t) to[bend left=20] node[fl] {2} (c);
\end{tikzpicture}}
\only<3>{%
\begin{tikzpicture}[scale=0.95]
\useasboundingbox (-0.45,-0.5) rectangle (5.85,3.3);
  \node[fn] (s) at (0,1.4) {$s$};
  \node[fn] (a) at (2.6,2.8) {$a$};
  \node[fn] (b) at (1.6,0) {$b$};
  \node[fn] (c) at (3.9,0) {$c$};
  \node[fn] (t) at (5.4,1.4) {$t$};
  \draw[res] (s) to[bend left=20] node[fl] {1} (a);
  \draw[resb] (a) to[bend left=20] node[fl] {2} (s);
  \draw[aug] (s) to[bend left=20] node[fl,text=kthgrass,font=\scriptsize\bfseries] {5} (b);
  \draw[resb] (b) to[bend left=20] node[fl] {1} (s);
  \draw[res] (b) -- node[fl] {1} (a);
  \draw[resb] (t) -- node[fl] {1} (a);
  \draw[res] (a) to[bend left=20] node[fl] {1} (c);
  \draw[resb] (c) to[bend left=20] node[fl] {1} (a);
  \draw[aug] (b) to[bend left=20] node[fl,text=kthgrass,font=\scriptsize\bfseries] {4} (c);
  \draw[resb] (c) to[bend left=20] node[fl] {1} (b);
  \draw[aug] (c) to[bend left=20] node[fl,text=kthgrass,font=\scriptsize\bfseries] {5} (t);
  \draw[resb] (t) to[bend left=20] node[fl] {2} (c);
\end{tikzpicture}}
\only<4>{%
\begin{tikzpicture}[scale=0.95]
\useasboundingbox (-0.45,-0.5) rectangle (5.85,3.3);
  \node[fn] (s) at (0,1.4) {$s$};
  \node[fn] (a) at (2.6,2.8) {$a$};
  \node[fn] (b) at (1.6,0) {$b$};
  \node[fn] (c) at (3.9,0) {$c$};
  \node[fn] (t) at (5.4,1.4) {$t$};
  \draw[res] (s) to[bend left=20] node[fl] {1} (a);
  \draw[resb] (a) to[bend left=20] node[fl] {2} (s);
  \draw[aug] (s) to[bend left=20] node[fl,text=kthgrass,font=\scriptsize\bfseries] {1} (b);
  \draw[resb] (b) to[bend left=20] node[fl] {5} (s);
  \draw[aug] (b) -- node[fl,text=kthgrass,font=\scriptsize\bfseries] {1} (a);
  \draw[resb] (t) -- node[fl] {1} (a);
  \draw[aug] (a) to[bend left=20] node[fl,text=kthgrass,font=\scriptsize\bfseries] {1} (c);
  \draw[resb] (c) to[bend left=20] node[fl] {1} (a);
  \draw[resb] (c) -- node[fl] {5} (b);
  \draw[aug] (c) to[bend left=20] node[fl,text=kthgrass,font=\scriptsize\bfseries] {1} (t);
  \draw[resb] (t) to[bend left=20] node[fl] {6} (c);
\end{tikzpicture}}
\only<5>{%
\begin{tikzpicture}[scale=0.95]
\useasboundingbox (-0.45,-0.5) rectangle (5.85,3.3);
  \node[fn] (s) at (0,1.4) {$s$};
  \node[fn] (a) at (2.6,2.8) {$a$};
  \node[fn] (b) at (1.6,0) {$b$};
  \node[fn] (c) at (3.9,0) {$c$};
  \node[fn] (t) at (5.4,1.4) {$t$};
  \begin{scope}[on background layer]
  \filldraw[plum!13,line width=0.8cm,line join=round] (0,1.4) -- (1.6,0) -- (2.6,2.8) -- cycle;
  \end{scope}
  \draw[res] (s) to[bend left=20] node[fl] {1} (a);
  \draw[resb] (a) to[bend left=20] node[fl] {2} (s);
  \draw[resb] (b) -- node[fl] {6} (s);
  \draw[resb] (a) -- node[fl] {1} (b);
  \draw[resb] (t) -- node[fl] {1} (a);
  \draw[resb] (c) -- node[fl] {2} (a);
  \draw[resb] (c) -- node[fl] {5} (b);
  \draw[resb] (t) -- node[fl] {7} (c);
\end{tikzpicture}}
\end{column}
\end{columns}
\vfill
\begin{tcolorbox}[enhanced,halign=left,colback=kthlight!50,colframe=kthsky,boxrule=0.4pt,arc=2pt,
  left=4pt,right=4pt,top=1pt,bottom=1pt,fontupper=\small,height=1.05cm,valign=center]
\only<1>{Inget flöde än: restgrafen är bara kapaciteterna. Välj en stig, vilken som helst: $s\to a\to c\to t$, flaskhals $\min(3,2,7)=2$.}\only<2>{$|f|=2$. Bakåtkanten $c\to a$ (orange) betyder ”ångra flöde på $a\to c$”. Stig $s\to b\to c\to a\to t$ använder den: $+1$.}\only<3>{$|f|=3$. Stig $s\to b\to c\to t$, flaskhals $\min(5,4,5)=4$.}\only<4>{$|f|=7$. Stig $s\to b\to a\to c\to t$, flaskhals 1.}\only<5>{$|f|=8$. Ingen stig $s\leadsto t$ i $G_f$. Nåbara från $s$: $\{s,a,b\}$ — ett minsnitt med kapacitet $1+2+5=8$.}
\end{tcolorbox}
\end{frame}


\begin{frame}{Repetition: maxflöde = minsnitt}
\begin{columns}[T]
\begin{column}{0.44\textwidth}
\centering
\input{gen/minsnitt}\par\smallskip
\footnotesize Två minsnitt, $C_1$ och $C_2$, båda 8.\\ Att de är två kommer att spela roll.
\end{column}
\begin{column}{0.53\textwidth}
\footnotesize
\renewcommand{\arraystretch}{1.1}
\begin{tabular}{@{}lll@{}}
\toprule
$S$ ($s$-sidan) & kanter $S\to T$ & kapacitet\\
\midrule
$\{s\}$ & $s{\to}a$, $s{\to}b$ & $3+6=9$\\
$\{s,a\}$ & $s{\to}b$, $a{\to}t$, $a{\to}c$ & $6+1+2=9$\\
$\{s,b\}$ & $s{\to}a$, $b{\to}a$, $b{\to}c$ & $3+1+5=9$\\
$\{s,c\}$ & $s{\to}a$, $s{\to}b$, $c{\to}t$ & $3+6+7=16$\\
\rowcolor{plum!12} $\{s,a,b\}$ & $a{\to}t$, $a{\to}c$, $b{\to}c$ & $1+2+5=\mathbf{8}$\enspace$C_1$\\
$\{s,a,c\}$ & $s{\to}b$, $a{\to}t$, $c{\to}t$ & $6+1+7=14$\\
$\{s,b,c\}$ & $s{\to}a$, $b{\to}a$, $c{\to}t$ & $3+1+7=11$\\
\rowcolor{plum!12} $\{s,a,b,c\}$ & $a{\to}t$, $c{\to}t$ & $1+7=\mathbf{8}$\enspace$C_2$\\
\bottomrule
\end{tabular}
\end{column}
\end{columns}
\vfill
\begin{missbox}
”Snittets kapacitet är summan av kanterna som korsar det.” Bara kanterna från $S$ till $T$
räknas. För $S=\{s,a\}$ korsar $b\to a$ snittet baklänges och räknas inte.
\end{missbox}
\end{frame}

\begin{frame}{Uppgift 2: Förändrat flöde}
\begin{infobox}
Anta att vi har en lösning $\Phi$ till flödesproblemet för en graf. Beskriv en effektiv algoritm
som hittar ett nytt maximalt flöde om kapaciteten längs en viss kant \textbf{a)}~ökar med en enhet,
\textbf{b)}~minskar med en enhet. Tidskomplexiteten ska vara linjär, alltså $O(|V|+|E|)$.
\end{infobox}
{\small\textbf{Vad består $\Phi$ av?} Två maximala flöden i samma nät:}\par
\begin{columns}[T]
\begin{column}{0.32\textwidth}
\centering
{\footnotesize\bfseries\color{kthnavy}$\Phi_1$, värde 8}\par
\begin{tikzpicture}[scale=0.66]
\useasboundingbox (-0.45,-0.5) rectangle (5.85,3.3);
  \node[fn] (s) at (0,1.4) {$s$};
  \node[fn] (a) at (2.6,2.8) {$a$};
  \node[fn] (b) at (1.6,0) {$b$};
  \node[fn] (c) at (3.9,0) {$c$};
  \node[fn] (t) at (5.4,1.4) {$t$};
  \draw[fe] (s) -- node[fl] {2/3} (a);
  \draw[sat] (s) -- node[fl] {6/6} (b);
  \draw[sat] (b) -- node[fl] {1/1} (a);
  \draw[sat] (a) -- node[fl] {1/1} (t);
  \draw[sat] (a) -- node[fl] {2/2} (c);
  \draw[sat] (b) -- node[fl] {5/5} (c);
  \draw[sat] (c) -- node[fl] {7/7} (t);
\end{tikzpicture}

\end{column}
\begin{column}{0.32\textwidth}
\centering
{\footnotesize\bfseries\color{kthnavy}$\Phi_2$, värde 8}\par
\input{gen/fstar}
\end{column}
\begin{column}{0.31\textwidth}
\footnotesize
Samma värde. Olika flöden: $s{\to}a$ 2 mot 3, $s{\to}b$ 6 mot 5, $b{\to}a$ 1 mot 0.\par\smallskip
$\Phi$ är flödet $f(e)$ \emph{på varje kant}. Med bara talet 8 finns inget att uppdatera — då
återstår att räkna om från början.
\end{column}
\end{columns}
\varstrip{Kontrast}{nätet och värdet 8}{flödet på kanterna}{att $\Phi$ är en funktion på kanterna, inte ett tal}
\end{frame}

% ---------------------------------------------------------------- linjär i vad
\begin{frame}{Linjär i vad?}
\begin{columns}[T]
\begin{column}{0.22\textwidth}
\centering
{\footnotesize\bfseries\color{kthnavy}Grafen}\par\smallskip
\begin{tikzpicture}[scale=0.95]
  \node[gv] (1) at (0,1.3) {1}; \node[gv] (2) at (1.3,1.3) {2};
  \node[gv] (3) at (0,0) {3};   \node[gv] (4) at (1.3,0) {4};
  \draw[ge] (1)--(2) (1)--(3) (2)--(3) (3)--(4);
\end{tikzpicture}\par
{\scriptsize $|V|=4$, $|E|=4$}
\end{column}
\begin{column}{0.36\textwidth}
\centering
{\footnotesize\bfseries\color{kthnavy}Grannlistor: $|V|+2|E|=12$ platser}\par\smallskip
\begin{tikzpicture}[cell/.style={draw=kthnavy,minimum width=4.4mm,minimum height=4.4mm,
    inner sep=0pt,font=\scriptsize},head/.style={cell,fill=kthlight},ent/.style={cell,fill=kthmint!70}]
  \foreach \v/\ns [count=\r] in {1/{2,3},2/{1,3},3/{1,2,4},4/{3}} {
    \node[head] at (0,-0.5*\r) {\v};
    \draw[->,kthnavy,line width=0.5pt] (0.25,-0.5*\r) -- (0.52,-0.5*\r);
    \foreach \n [count=\k] in \ns { \node[ent] at (0.3+0.44*\k,-0.5*\r) {\n}; }
  }
\end{tikzpicture}
\end{column}
\begin{column}{0.36\textwidth}
\centering
{\footnotesize\bfseries\color{kthnavy}Grannmatris: $|V|^2=16$ platser}\par\smallskip
\begin{tikzpicture}[cell/.style={draw=kthnavy,minimum width=4.4mm,minimum height=4.4mm,
    inner sep=0pt,font=\scriptsize},c0/.style={fill=white,text=black!40},c1/.style={fill=kthmint!70}]
  \foreach \row [count=\i] in {{0,1,1,0},{1,0,1,0},{1,1,0,1},{0,0,1,0}} {
    \node[font=\scriptsize,text=black!50] at (0,-0.44*\i) {\i};
    \node[font=\scriptsize,text=black!50] at (0.44*\i,0) {\i};
    \foreach \x [count=\j] in \row { \node[cell,c\x] at (0.44*\j,-0.44*\i) {\x}; }
  }
\end{tikzpicture}
\end{column}
\end{columns}
\smallskip
\footnotesize
\begin{itemize}\setlength{\itemsep}{1pt}
  \item \textbf{Linjär tid} = högst en konstant gånger \emph{indatas storlek}. Man hinner läsa
        indata ett fast antal gånger, inte mer.
  \item BFS/DFS läser varje plats en gång: $O(|V|+|E|)$ med grannlistor, $O(|V|^2)$ med
        grannmatris. \emph{Båda är linjära} — i storleken på sin representation.
  \item $O(|V|+|E|)$ är inte ”linjärt i $|V|$”. I en fullständig graf är
        $|E|=\binom{|V|}{2}$, och då är $O(|V|+|E|)=O(|V|^2)$.
\end{itemize}
\varstrip{Separation}{grafen och algoritmen (BFS)}{representationen}{att ”linjär” mäts i indatas storlek}
\end{frame}

\begin{frame}{Linjär i vad? Dagens uppgifter}
\footnotesize
\renewcommand{\arraystretch}{1.2}
\begin{tabularx}{\textwidth}{@{}p{0.17\textwidth}p{0.3\textwidth}p{0.15\textwidth}X@{}}
\toprule
\textbf{Uppgift} & \textbf{Indatas storlek} & \textbf{”Linjärt”} & \textbf{Kommentar}\\
\midrule
2 Förändrat flöde & $|V|+|E|$ (grafen och $\Phi$) & $O(|V|+|E|)$ & ett par grafsökningar\\
3 Eulercykel & $|V|+|E|$ & $O(|V|+|E|)$ & utdata har $|E|+1$ hörn\\
4 Julklappar & $n$ listor à högst $m$ saker & $O(n+nm)$ & ett varv av Ford--Fulkerson; det blir $n$ varv\\
5 Sökning & $n$ sorterade tal & — & $\Theta(\log n)$: läser inte ens all indata\\
6 Lampor & $n$ lampor & $\Theta(n)$ byten & uppgiften vill ha exakt antal: $2n-1$\\
\bottomrule
\end{tabularx}
\medskip
\begin{missbox}
”Eulercykeln hittas i $O(n)$.” (Så står det i äldre övningsbilder.) Om $n=|V|$ är det fel: en graf
med $|V|$ hörn kan ha $\Theta(|V|^2)$ kanter, och alla ska med i cykeln.
\end{missbox}
\varstrip{Separation}{ordet ”linjär”}{storleksmåttet: $|V|+|E|$, $|V|^2$, $n$, antal byten}{att ”linjär” alltid är linjär \emph{i något}}
\end{frame}

{\kaninbg
\begin{frame}{\kh Linjär i ett tal: pseudopolynomisk tid}
\begin{columns}[T]
\begin{column}{0.38\textwidth}
\centering
\begin{tikzpicture}[scale=1.05]
  \node[fn] (s) at (0,1.1) {$s$};   \node[fn] (a) at (1.7,2.2) {$a$};
  \node[fn] (b) at (1.7,0) {$b$};   \node[fn] (t) at (3.4,1.1) {$t$};
  \draw[fe] (s) -- node[fl] {$C$} (a);  \draw[fe] (s) -- node[fl] {$C$} (b);
  \draw[fe] (a) -- node[fl] {$C$} (t);  \draw[fe] (b) -- node[fl] {$C$} (t);
  \draw[sat] (a) -- node[fl] {1} (b);
\end{tikzpicture}\par\smallskip
{\footnotesize $C=1\,000\,000\approx2^{20}$}
\end{column}
\begin{column}{0.59\textwidth}
\small
\begin{itemize}\setlength{\itemsep}{2pt}
  \item Ford--Fulkerson väljer stigar godtyckligt. Med otur: varannan stig $s\to a\to b\to t$,
        varannan $s\to b\to a\to t$ (baklänges på $a\to b$).
  \item Varje stig ger $+1$. Maxflödet är $2C$: två miljoner varv.
  \item Indata: fem kanter, tal på cirka 20 bitar. Linjärt i \emph{talet} $C$, exponentiellt
        i antalet \emph{bitar}.
  \item Edmonds--Karp (kortaste stigen först) behöver två varv här, och $O(|V|\,|E|^2)$ i allmänhet.
\end{itemize}
\end{column}
\end{columns}
\vfill
\begin{kaninbox}[fontupper=\footnotesize]
I uppgift 4 är alla kapaciteter 1 och $|f^*|\le n$. Där är $O(|E|\cdot|f^*|)$ helt ofarligt.
\end{kaninbox}
\varstrip{Kontrast}{nätet}{stigvalet: godtyckligt eller kortast}{linjärt i ett tal $\neq$ linjärt i indatas längd}
\end{frame}}

% ---------------------------------------------------------------- 2a
\begin{frame}{2a: kapaciteten på $(u,v)$ ökar med 1}
\begin{columns}[T]
\begin{column}{0.5\textwidth}
\small
\begin{itemize}\setlength{\itemsep}{3pt}
  \item $\Phi$ är fortfarande ett tillåtet flöde: kapaciteten blev ju bara större.
  \item I restflödesgrafen $G_\Phi$ ändras en enda sak: restkapaciteten $u\to v$ ökar med 1.
  \item Kör \textbf{en} iteration av Ford--Fulkerson: sök en stig $s\leadsto t$ i $G_\Phi$.
  \item Finns en stig: öka flödet med 1 längs den. Annars är $\Phi$ fortfarande maximalt.
\end{itemize}
\end{column}
\begin{column}{0.47\textwidth}
{\small\bfseries\color{kthnavy}ÖkaKapacitet$(G,\Phi,(u,v))$}
\footnotesize
\begin{algorithmic}[1]
\State $c(u,v)\gets c(u,v)+1$
\State bygg $G_\Phi$ \Comment{$O(|V|+|E|)$}
\State $P\gets\textsc{BFS}(G_\Phi,s,t)$ \Comment{$O(|V|+|E|)$}
\If{$P$ finns}
  \State öka $\Phi$ med 1 längs $P$ \Comment{$O(|V|)$}
\EndIf
\State \Return $\Phi$
\end{algorithmic}
\end{column}
\end{columns}
\vfill
\begin{columns}[T]
\begin{column}{0.49\textwidth}
\begin{ahabox}[fontupper=\footnotesize]
\textbf{Genväg:} var $(u,v)$ inte mättad händer ingenting. Minsnittet som Ford--Fulkerson hittade
(de nåbara hörnen) har bara mättade kanter ut, så dess kapacitet är oförändrad.
\end{ahabox}
\end{column}
\begin{column}{0.49\textwidth}
\begin{missbox}[fontupper=\footnotesize]
”Kör Ford--Fulkerson igen.” Från början tar det $O(|E|\cdot|f^*|)$ och kastar bort $\Phi$.
\end{missbox}
\end{column}
\end{columns}
\end{frame}

\begin{frame}{2a, exempel: $a\to t$ får kapacitet 2}
\begin{columns}[T]
\begin{column}{0.32\textwidth}
\centering
{\footnotesize\bfseries\color{kthnavy}$\Phi_1$ med ny kapacitet}\par
\input{gen/2a-net}
\end{column}
\begin{column}{0.32\textwidth}
\centering
{\footnotesize\bfseries\color{kthnavy}$G_\Phi$: stigen $s\to a\to t$}\par
\input{gen/2a-res}
\end{column}
\begin{column}{0.32\textwidth}
\centering
{\footnotesize\bfseries\color{kthnavy}Nytt flöde, värde 9}\par
\input{gen/2a-ny}
\end{column}
\end{columns}
\vfill
\small
Restkanten $a\to t$ (kapacitet 1) är ny, och $s\to a$ hade redan 1 kvar. BFS hittar
$s\to a\to t$. Kanten $a\to t$ ligger i båda minsnitten $C_1$ och $C_2$, så båda växte till 9.
\varstrip{Generalisering}{metoden}{exemplet}{en sökning, en enhet}
\end{frame}

\begin{frame}{2a, exempel: $a\to c$ får kapacitet 3}
\begin{columns}[T]
\begin{column}{0.32\textwidth}
\centering
{\footnotesize\bfseries\color{kthnavy}$\Phi_1$ med ny kapacitet}\par
\input{gen/2a-ac-net}
\end{column}
\begin{column}{0.32\textwidth}
\centering
{\footnotesize\bfseries\color{kthnavy}$G_\Phi$: nåbart från $s$}\par
\input{gen/2a-res-ac}
\end{column}
\begin{column}{0.32\textwidth}
\small
Nu nås $c$ via $s\to a\to c$ — men $c\to t$ är mättad. $t$ nås inte.\par\medskip
$\Phi_1$ är fortfarande maximalt: 8. Minsnittet $C_2=\{s,a,b,c\}$ innehåller inte $a\to c$
och har kvar kapaciteten 8.
\end{column}
\end{columns}
\vfill
\begin{missbox}
”Kanten var mättad, alltså hjälper det att öka den.” Mättad är nödvändigt men inte tillräckligt.
\end{missbox}
\varstrip{Kontrast}{nätet och $\Phi_1$}{vilken mättad kant som ökas}{att en ökning kan vara verkningslös}
\end{frame}

{\kaninbg
\begin{frame}{\kh Varför räcker en enda sökning?}
\begin{columns}[T]
\begin{column}{0.56\textwidth}
\small
\begin{itemize}\setlength{\itemsep}{3pt}
  \item Ett snitts kapacitet är summan av kanterna $S\to T$. När $c(u,v)$ ökar med 1 växer
        varje snitt med 1 (om $u\in S$, $v\in T$) eller med 0.
  \item Maxflöde = minsnitt, så det nya maxflödet är högst $|\Phi|+1$.
  \item Allt är heltal: antingen samma eller exakt $+1$. En förbättrande stig, sedan är vi klara.
  \item Ingen stig? Då är $\Phi$ maximalt. Det är Ford--Fulkersons eget stoppvillkor.
\end{itemize}
\end{column}
\begin{column}{0.41\textwidth}
\centering
\begin{tikzpicture}[scale=0.8]
  \node[fn] (s) at (0,1.4) {$s$};
  \node[fn] (a) at (2.6,2.8) {$a$};
  \node[fn] (b) at (1.6,0) {$b$};
  \node[fn] (c) at (3.9,0) {$c$};
  \node[fn] (t) at (5.4,1.4) {$t$};
  \draw[fe] (s) -- node[fl] {2/3} (a);
  \draw[sat] (s) -- node[fl] {6/6} (b);
  \draw[sat] (b) -- node[fl] {1/1} (a);
  \draw[sat] (a) -- node[fl] {1/1} (t);
  \draw[sat] (a) -- node[fl] {2/2} (c);
  \draw[sat] (b) -- node[fl] {5/5} (c);
  \draw[sat] (c) -- node[fl] {7/7} (t);
  \draw[cutline] plot[smooth,tension=0.7] coordinates {(3.3,3.35) (2.95,2.2) (2.9,1.3) (3.3,0.1) (3.35,-0.5)};
  \node[below,font=\scriptsize\bfseries,text=plum] at (3.35,-0.5) {$C_1$};
  \draw[cutline] (4.98,3.3) -- (4.98,-0.5) node[below,font=\scriptsize\bfseries,text=plum] {$C_2$};
\end{tikzpicture}

\end{column}
\end{columns}
\vfill
\begin{kaninbox}[fontupper=\footnotesize]
\textbf{Och om kapaciteten ökar med $k$?} Då kan maxflödet öka med upp till $k$, och det kan behövas
$k$ sökningar: $O(k(|V|+|E|))$. Linjärt bara om $k$ är en konstant.
\end{kaninbox}
\end{frame}}

\begin{frame}{Samma ändring, olika kant: när hjälper $+1$?}
\begin{columns}[T]
\begin{column}{0.4\textwidth}
\centering
\begin{tikzpicture}[scale=0.8]
  \node[fn] (s) at (0,1.4) {$s$};
  \node[fn] (a) at (2.6,2.8) {$a$};
  \node[fn] (b) at (1.6,0) {$b$};
  \node[fn] (c) at (3.9,0) {$c$};
  \node[fn] (t) at (5.4,1.4) {$t$};
  \draw[fe] (s) -- node[fl] {2/3} (a);
  \draw[sat] (s) -- node[fl] {6/6} (b);
  \draw[sat] (b) -- node[fl] {1/1} (a);
  \draw[sat] (a) -- node[fl] {1/1} (t);
  \draw[sat] (a) -- node[fl] {2/2} (c);
  \draw[sat] (b) -- node[fl] {5/5} (c);
  \draw[sat] (c) -- node[fl] {7/7} (t);
  \draw[cutline] plot[smooth,tension=0.7] coordinates {(3.3,3.35) (2.95,2.2) (2.9,1.3) (3.3,0.1) (3.35,-0.5)};
  \node[below,font=\scriptsize\bfseries,text=plum] at (3.35,-0.5) {$C_1$};
  \draw[cutline] (4.98,3.3) -- (4.98,-0.5) node[below,font=\scriptsize\bfseries,text=plum] {$C_2$};
\end{tikzpicture}
\par
\begin{ahabox}[fontupper=\footnotesize]
Maxflödet ökar precis när kanten ligger i \textbf{alla} minsnitt. Ett enda minsnitt utan
kanten håller taket kvar.
\end{ahabox}
\end{column}
\begin{column}{0.57\textwidth}
\footnotesize
\renewcommand{\arraystretch}{1.1}
\begin{tabular}{@{}lccccc@{}}
\toprule
Kant & $\Phi_1$ & mättad & i $C_1$ & i $C_2$ & $+1$ ger\\
\midrule
\rowcolor{kthmint!50} $a{\to}t$ & 1/1 & \ja & \ja & \ja & \textbf{9}\\
$a{\to}c$ & 2/2 & \ja & \ja & & 8\\
$b{\to}c$ & 5/5 & \ja & \ja & & 8\\
$c{\to}t$ & 7/7 & \ja & & \ja & 8\\
$s{\to}b$ & 6/6 & \ja & & & 8\\
$b{\to}a$ & 1/1 & \ja & & & 8\\
$s{\to}a$ & 2/3 & & & & 8\\
\bottomrule
\end{tabular}
\par\smallskip
\begin{missbox}[fontupper=\footnotesize]
”Mättad kant $\Rightarrow$ ökningen hjälper.” Sex av sju kanter är mättade. En hjälper.
\end{missbox}
\end{column}
\end{columns}
\varstrip{Separation}{nätet, $\Phi_1$ och ändringen $+1$}{kanten}{att minsnitten bestämmer, inte mättnaden}
\end{frame}

% ---------------------------------------------------------------- 2b
\begin{frame}{2b: kapaciteten på $(u,v)$ minskar med 1}
\begin{columns}[T]
\begin{column}{0.49\textwidth}
\small
\textbf{Fall 1:} $(u,v)$ är inte mättad. $\Phi$ är fortfarande tillåtet, och maxflödet kan inte
öka när en kapacitet minskar. Klart.\par\smallskip
\textbf{Annars:} minska $\Phi(u,v)$ med 1. Nu kommer det in en enhet mer till $u$ än som går ut,
och en enhet mindre till $v$.\par\smallskip
\textbf{Fall 2:} finns en stig $u\leadsto v$ i $G_\Phi$? Led enheten den vägen. Värdet är
oförändrat.\par\smallskip
\textbf{Fall 3:} annars: skicka tillbaka enheten, $u\leadsto s$, och ta bort en enhet från $v$
till $t$ genom att gå $t\leadsto v$. Värdet minskar med 1.
\end{column}
\begin{column}{0.49\textwidth}
{\small\bfseries\color{kthnavy}MinskaKapacitet$(G,\Phi,(u,v))$}
\footnotesize
\begin{algorithmic}[1]
\State $c(u,v)\gets c(u,v)-1$
\If{$\Phi(u,v)\le c(u,v)$} \State \Return $\Phi$ \Comment{fall 1}
\EndIf
\State $\Phi(u,v)\gets\Phi(u,v)-1$ \Comment{obalans i $u$, $v$}
\If{stig $P\colon u\leadsto v$ finns i $G_\Phi$}
  \State öka $\Phi$ med 1 längs $P$ \Comment{fall 2}
\Else
  \State öka $\Phi$ med 1 längs en stig $u\leadsto s$ \Comment{fall 3}
  \State öka $\Phi$ med 1 längs en stig $t\leadsto v$
\EndIf
\State \Return $\Phi$
\end{algorithmic}
\smallskip
{\footnotesize Högst tre grafsökningar: $O(|V|+|E|)$.}
\end{column}
\end{columns}
\vfill
\begin{missbox}
”Minska flödet på kanten med 1, klart.” Då kommer mer in i $u$ än som går ut. Det är inte längre
ett flöde.
\end{missbox}
\end{frame}

\begin{frame}{2b, exempel (fall 2): $b\to a$ får kapacitet 0}
\begin{columns}[T]
\begin{column}{0.32\textwidth}
\centering
{\footnotesize\bfseries\color{kthnavy}Obalans efter minskningen}\par
\input{gen/2b-ba-net}
\end{column}
\begin{column}{0.32\textwidth}
\centering
{\footnotesize\bfseries\color{kthnavy}$G_\Phi$: stigen $b\to s\to a$}\par
\begin{tikzpicture}[scale=0.72]
\useasboundingbox (-0.45,-0.5) rectangle (5.85,3.3);
  \node[fn] (s) at (0,1.4) {$s$};
  \node[fn] (a) at (2.6,2.8) {$a$};
  \node[fn] (b) at (1.6,0) {$b$};
  \node[fn] (c) at (3.9,0) {$c$};
  \node[fn] (t) at (5.4,1.4) {$t$};
  \draw[aug] (s) to[bend left=20] node[fl,text=kthgrass,font=\scriptsize\bfseries] {1} (a);
  \draw[resb] (a) to[bend left=20] node[fl] {2} (s);
  \draw[aug] (b) -- node[fl,text=kthgrass,font=\scriptsize\bfseries] {6} (s);
  \draw[resb] (t) -- node[fl] {1} (a);
  \draw[resb] (c) -- node[fl] {2} (a);
  \draw[resb] (c) -- node[fl] {5} (b);
  \draw[resb] (t) -- node[fl] {7} (c);
\end{tikzpicture}

\end{column}
\begin{column}{0.32\textwidth}
\centering
{\footnotesize\bfseries\color{kthnavy}Nytt flöde, värde 8}\par
\input{gen/2b-ba-ny}
\end{column}
\end{columns}
\vfill
\small
$b$ har en enhet för mycket, $a$ en för lite. Stigen $b\to s\to a$ går baklänges på $s\to b$ och
framlänges på $s\to a$: enheten tar vägen $s\to a$ i stället. Resultatet är $\Phi_2$ — det andra
maxflödet från bilden om vad $\Phi$ är.
\varstrip{Kontrast}{nätet, $\Phi_1$ och minskningen}{kanten (här: ingen i något minsnitt)}{omledning bevarar värdet}
\end{frame}

\begin{frame}{2b, exempel (fall 3): $a\to c$ får kapacitet 1}
\begin{columns}[T]
\begin{column}{0.32\textwidth}
\centering
{\footnotesize\bfseries\color{kthnavy}Obalans efter minskningen}\par
\input{gen/2b-ac-net}
\end{column}
\begin{column}{0.32\textwidth}
\centering
{\footnotesize\bfseries\color{kthnavy}$G_\Phi$: nåbart från $a$}\par
\begin{tikzpicture}[scale=0.72]
\useasboundingbox (-0.45,-0.5) rectangle (5.85,3.3);
  \node[fn] (s) at (0,1.4) {$s$};
  \node[fn] (a) at (2.6,2.8) {$a$};
  \node[fn] (b) at (1.6,0) {$b$};
  \node[fn] (c) at (3.9,0) {$c$};
  \node[fn] (t) at (5.4,1.4) {$t$};
  \begin{scope}[on background layer]
  \filldraw[plum!13,line width=0.8cm,line join=round] (0,1.4) -- (1.6,0) -- (2.6,2.8) -- cycle;
  \end{scope}
  \draw[res] (s) to[bend left=20] node[fl] {1} (a);
  \draw[aug] (a) to[bend left=20] node[fl,text=kthgrass,font=\scriptsize\bfseries] {2} (s);
  \draw[resb] (b) -- node[fl] {6} (s);
  \draw[resb] (a) -- node[fl] {1} (b);
  \draw[resb] (t) -- node[fl] {1} (a);
  \draw[resb] (c) -- node[fl] {1} (a);
  \draw[resb] (c) -- node[fl] {5} (b);
  \draw[aug] (t) -- node[fl,text=kthgrass,font=\scriptsize\bfseries] {7} (c);
\end{tikzpicture}

\end{column}
\begin{column}{0.32\textwidth}
\centering
{\footnotesize\bfseries\color{kthnavy}Nytt flöde, värde 7}\par
\input{gen/2b-ac-ny}
\end{column}
\end{columns}
\vfill
\small
Från $a$ nås bara $\{a,s,b\}$, inte $c$. Alltså fall 3: skicka tillbaka längs $a\to s$ och ta
bort en enhet längs $t\to c$ (båda baklänges i restgrafen, gröna). Kanten $a\to c$ ligger i
minsnittet $C_1$, som nu har kapacitet 7.
\varstrip{Kontrast}{nätet, $\Phi_1$ och minskningen}{kanten (här: i minsnittet $C_1$)}{när värdet måste sjunka}
\end{frame}

{\kaninbg
\begin{frame}{\kh Varför finns stigarna $u\leadsto s$ och $t\leadsto v$?}
\begin{columns}[T]
\begin{column}{0.57\textwidth}
\small
\begin{itemize}\setlength{\itemsep}{3pt}
  \item \textbf{Flödesuppdelning:} varje flöde är en summa av $s$--$t$-stigar och cykler med
        positivt flöde. $\Phi_1$: $s{\to}b{\to}c{\to}t$ (5), $s{\to}a{\to}c{\to}t$ (2),
        $s{\to}b{\to}a{\to}t$ (1).
  \item $(u,v)$ var mättad, så $\Phi(u,v)\ge1$: någon stig eller cykel går genom $(u,v)$.
  \item \textbf{Stig} $s\leadsto u\to v\leadsto t$: flödet på delarna ger bakåtkanter
        $u\leadsto s$ och $t\leadsto v$ i $G_\Phi$.
  \item \textbf{Cykel} genom $(u,v)$: resten av cykeln, baklänges, är en stig $u\leadsto v$.
        Då var vi i fall 2.
\end{itemize}
\end{column}
\begin{column}{0.4\textwidth}
\centering
{\footnotesize\bfseries\color{kthgrass}$s\to a\to c\to t$ bär 2 enheter}\par
\begin{tikzpicture}[scale=0.8]
\useasboundingbox (-0.45,-0.5) rectangle (5.85,3.3);
  \node[fn] (s) at (0,1.4) {$s$};
  \node[fn] (a) at (2.6,2.8) {$a$};
  \node[fn] (b) at (1.6,0) {$b$};
  \node[fn] (c) at (3.9,0) {$c$};
  \node[fn] (t) at (5.4,1.4) {$t$};
  \draw[aug] (s) -- node[fl] {2/3} (a);
  \draw[fe] (s) -- node[fl] {6/6} (b);
  \draw[fe] (b) -- node[fl] {1/1} (a);
  \draw[fe] (a) -- node[fl] {1/1} (t);
  \draw[aug] (a) -- node[fl] {2/2} (c);
  \draw[fe] (b) -- node[fl] {5/5} (c);
  \draw[aug] (c) -- node[fl] {7/7} (t);
\end{tikzpicture}

\end{column}
\end{columns}
\vfill
\begin{kaninbox}[fontupper=\footnotesize]
Gör sökningarna en i taget: först $u\leadsto s$, uppdatera $\Phi$, sedan $t\leadsto v$ i den nya
restgrafen. Samma uppdelningsargument på det uppdaterade flödet visar att den andra stigen
fortfarande finns.
\end{kaninbox}
\end{frame}}

{\kaninbg
\begin{frame}{\kh Är flödet i fall 3 säkert maximalt?}
\begin{columns}[T]
\begin{column}{0.58\textwidth}
\small
\footnotesize
\begin{itemize}\setlength{\itemsep}{2pt}
  \item Låt $A$ = hörnen som nås från $u$ i $G_\Phi$. Fall 3 betyder $v\notin A$.
  \item $s\in A$, för stigen $u\leadsto s$ finns. $t\notin A$, annars fanns $u\leadsto t\leadsto v$.
  \item Ingen restkant lämnar $A$: kanterna ut ur $A$ är mättade, kanterna in är tomma.
  \item Flödet över snittet är då lika med dess kapacitet, och det är $|\Phi|-1$: enheten som
        blev över i $u$ stannar i $A$.
  \item Tillbakatrycken sker helt inom $A$ respektive helt utanför. Nya flödet har värdet
        $|\Phi|-1$ = ett snitts kapacitet. Alltså maximalt.
\end{itemize}
\end{column}
\begin{column}{0.39\textwidth}
\centering
\begin{tikzpicture}[scale=0.8]
\useasboundingbox (-0.45,-0.5) rectangle (5.85,3.3);
  \node[fn] (s) at (0,1.4) {$s$};
  \node[fn] (a) at (2.6,2.8) {$a$};
  \node[fn] (b) at (1.6,0) {$b$};
  \node[fn] (c) at (3.9,0) {$c$};
  \node[fn] (t) at (5.4,1.4) {$t$};
  \begin{scope}[on background layer]
  \filldraw[plum!13,line width=0.8cm,line join=round] (0,1.4) -- (1.6,0) -- (2.6,2.8) -- cycle;
  \end{scope}
  \node[font=\small\bfseries,text=plum] at (0.35,0.1) {$A$};
  \draw[res] (s) to[bend left=20] node[fl] {1} (a);
  \draw[resb] (a) to[bend left=20] node[fl] {2} (s);
  \draw[resb] (b) -- node[fl] {6} (s);
  \draw[resb] (a) -- node[fl] {1} (b);
  \draw[resb] (t) -- node[fl] {1} (a);
  \draw[resb] (c) -- node[fl] {1} (a);
  \draw[resb] (c) -- node[fl] {5} (b);
  \draw[resb] (t) -- node[fl] {7} (c);
  \draw[cutline] plot[smooth,tension=0.7] coordinates {(3.3,3.35) (2.95,2.2) (2.9,1.3) (3.3,0.1) (3.35,-0.5)};
\end{tikzpicture}
\par
{\footnotesize Här: $A=\{s,a,b\}$, och $1+1+5=7$.}
\end{column}
\end{columns}
\vfill
\begin{kaninbox}[fontupper=\footnotesize]
Ingen extra kontroll behövs. (En extra BFS skulle inte förstöra linjäriteten. Den är bara onödig.)
\end{kaninbox}
\end{frame}}

{\kaninbg
\begin{frame}{\kh Alternativ lösning: b) via a)}
\small
Om $(u,v)$ är mättad:
\begin{enumerate}\setlength{\itemsep}{3pt}
  \item \textbf{Ta bort en enhet flöde genom $(u,v)$:} minska $\Phi(u,v)$ och tryck tillbaka
        längs $u\leadsto s$ och $t\leadsto v$. Nu är $\Phi$ tillåtet med den nya kapaciteten och
        har värdet $|\Phi|-1$.
  \item \textbf{Gör som i a):} en sökning $s\leadsto t$. Maxflödet är nu $|\Phi|-1$ eller
        $|\Phi|$, så en förbättrande stig räcker.
\end{enumerate}
\medskip
Exempel $b\to a$: steg 1 tar bort enheten på $s\to b\to a\to t$ (värde 7), steg 2 hittar
$s\to a\to t$ (värde 8). Samma slutresultat, $\Phi_2$.
\vfill
\begin{kaninbox}[fontupper=\footnotesize]
Tre sökningar i stället för en eller två — fortfarande $O(|V|+|E|)$. Vinsten är att
korrekthetsbeviset blir kortare: efter steg 1 är vi exakt i situationen i a).
\end{kaninbox}
\varstrip{Generalisering}{metoden: restgraf och en sökning}{problemet: öka eller minska}{att b) är ”ta bort en enhet” följt av a)}
\end{frame}}

% ---------------------------------------------------------------- fusion
\begin{frame}{Öka och minska, kant för kant}
\begin{columns}[T]
\begin{column}{0.58\textwidth}
\footnotesize
\renewcommand{\arraystretch}{1.02}
\begin{tabular}{@{}lcccccl@{}}
\toprule
Kant & $\Phi_1$ & i $C_1$ & i $C_2$ & $+1$ & $-1$ & fall vid $-1$\\
\midrule
$a{\to}t$ & 1/1 & \ja & \ja & \textbf{9} & \textbf{7} & 3\\
$a{\to}c$ & 2/2 & \ja & & 8 & \textbf{7} & 3\\
$b{\to}c$ & 5/5 & \ja & & 8 & \textbf{7} & 3\\
$c{\to}t$ & 7/7 & & \ja & 8 & \textbf{7} & 3\\
$s{\to}b$ & 6/6 & & & 8 & 8 & 2: $s{\to}a{\to}b$\\
$b{\to}a$ & 1/1 & & & 8 & 8 & 2: $b{\to}s{\to}a$\\
$s{\to}a$ & 2/3 & & & 8 & 8 & 1\\
\bottomrule
\end{tabular}
\end{column}
\begin{column}{0.39\textwidth}
\begin{ahabox}[fontupper=\footnotesize]
$+1$ ökar maxflödet $\iff$ kanten ligger i \textbf{alla} minsnitt.\par\smallskip
$-1$ minskar maxflödet $\iff$ kanten ligger i \textbf{något} minsnitt.
\end{ahabox}
\end{column}
\end{columns}
\vfill
\begin{columns}[T]
\begin{column}{0.49\textwidth}
\begin{missbox}[fontupper=\footnotesize]
”Ökning och minskning är symmetriska.” $a\to c$ kan inte öka maxflödet, men kan minska det.
\end{missbox}
\end{column}
\begin{column}{0.49\textwidth}
\begin{missbox}[fontupper=\footnotesize]
”Mättad är en egenskap hos kanten.” Hos kanten \emph{och} $\Phi$: $s\to b$ är mättad i $\Phi_1$
(6/6) men inte i $\Phi_2$ (5/6).
\end{missbox}
\end{column}
\end{columns}
\varstrip{Fusion}{nätet}{kanten och riktningen samtidigt}{alla mot något minsnitt; mättnad beror på $\Phi$}
\end{frame}

\begin{frame}{Vanliga fel i uppgift 2}
\begin{columns}[T]
\begin{column}{0.49\textwidth}
\begin{missbox}[fontupper=\footnotesize]
”$\Phi$ är ett tal, 8.” Det är flödet på varje kant. Med bara talet måste man räkna om från början.
\end{missbox}
\begin{missbox}[fontupper=\footnotesize]
”Öka flödet på $(u,v)$ med 1.” Då bryts flödesbevarandet i $u$ och $v$, och enheten kanske inte
ens kan komma fram till $t$.
\end{missbox}
\end{column}
\begin{column}{0.49\textwidth}
\begin{missbox}[fontupper=\footnotesize]
”Minska flödet på $(u,v)$ med 1, klart.” Samma fel åt andra hållet. Obalansen måste ledas om
(fall 2) eller skickas tillbaka (fall 3).
\end{missbox}
\begin{missbox}[fontupper=\footnotesize]
”Linjär tid, alltså $O(|V|)$.” Linjärt i indatas storlek, här $|V|+|E|$. Se ”Linjär i vad?”.
\end{missbox}
\end{column}
\end{columns}
\vfill
\begin{ahabox}[fontupper=\footnotesize]
Gemensam nämnare för de rätta lösningarna: ändra i restgrafen, gör en grafsökning, och låt
maxflöde = minsnitt tala om att det räcker.
\end{ahabox}
\varstrip{Kontrast}{uppgiften}{lösningsidén, rätt och nästan rätt}{vad som gör de rätta rätt}
\end{frame}
